\subsection{Reproducibility Checklist and Metadata}\label{app:reproducibility_checklist}

This appendix provides a comprehensive checklist and metadata snapshot to facilitate independent replication of the experimental results reported in \cref{sec:results}. All information documented here corresponds to the configuration and environment used for the final benchmark runs that produced the data analyzed in \cref{sec:statistical_analysis,sec:cost_modelling,sec:decision}.

\subsubsection{Software Environment}

\begin{table}[h!]
\centering
\caption{Software environment and dependency versions for reproducible execution.\label{tab:software_environment}}
\begin{tabularx}{\textwidth}{@{}l X@{}}
\toprule
\textbf{Component} & \textbf{Version / Details} \\
\midrule
Operating System & Ubuntu 22.04.3 LTS (Jammy Jellyfish), kernel 5.15.0-88 \\
Python Runtime & 3.10.12 (default, Jul 5 2023, 18:54:27) [GCC 11.2.0] \\
\texttt{pip} & 23.2.1 \\
\texttt{conda} & 23.7.4 (environment isolation for dependency management) \\
Harness Git Repository & \texttt{https://github.com/ksawesome/llm-harness} \\
Harness Git Commit SHA & \texttt{7ec2e576af5d6b8887d99f5b6bd6aaef4e4737f7} \\
Configuration Files Hash & SHA-256: \texttt{8f79cfa47c345eefd53d2595b813958a3ca8233da7d82d66d93bfb54dbe2d8d0} (models.json) \\
\bottomrule
\end{tabularx}
\end{table}

\paragraph{Core Python dependencies.} All libraries are pinned to exact versions in \texttt{requirements.txt} (excerpt shown in \cref{tab:python_dependencies}). Full dependency tree available via \texttt{pip freeze > requirements-frozen.txt} in the repository artifacts.

\begin{table}[h!]
\centering
\caption{Primary Python dependencies with exact pinned versions.\label{tab:python_dependencies}}
\begin{tabularx}{\textwidth}{@{}l l X@{}}
\toprule
\textbf{Package} & \textbf{Version} & \textbf{Purpose} \\
\midrule
\texttt{openai} & 1.14.3 & OpenAI GPT-4o Mini API client \\
\texttt{anthropic} & 0.21.3 & Anthropic Claude 3 Sonnet API client \\
\texttt{google-generativeai} & 0.3.2 & Google Gemini 2.5 Flash SDK \\
\texttt{cohere} & 4.38 & Cohere Command R 08-2024 API client \\
\texttt{huggingface-hub} & 0.20.1 & Hugging Face Inference API (Llama 4 Maverick) \\
\texttt{pandas} & 2.2.0 & Data manipulation and CSV export \\
\texttt{numpy} & 1.26.3 & Numerical computations and statistical proxies \\
\texttt{scikit-learn} & 1.4.0 & Statistical analysis and effect size computation \\
\texttt{scipy} & 1.12.0 & Hypothesis testing and confidence intervals \\
\texttt{flask} & 3.0.1 & Web dashboard server \\
\texttt{plotly} & 5.18.0 & Interactive visualizations \\
\texttt{pytest} & 7.4.0 & Unit and integration testing \\
\texttt{pytest-cov} & 4.1.0 & Code coverage measurement \\
\texttt{tiktoken} & 0.6.0 & OpenAI-compatible tokenization for cost estimation \\
\texttt{jsonschema} & 4.20.0 & Configuration validation \\
\texttt{python-dotenv} & 1.0.0 & Environment variable management \\
\bottomrule
\end{tabularx}
\end{table}

\subsubsection{Execution Context and Hardware}

\begin{table}[h!]
\centering
\caption{Execution hardware and network configuration.\label{tab:execution_context}}
\begin{tabularx}{\textwidth}{@{}l X@{}}
\toprule
\textbf{Parameter} & \textbf{Value / Details} \\
\midrule
\multicolumn{2}{l}{\textit{Hardware Specifications}} \\
Processor & AMD Ryzen 7 7840HS w/ Radeon 780M Graphics (8 cores, 16 threads) \\
Memory & 16 GB DDR5 RAM \\
Storage & 512 GB Samsung NVMe SSD \\
\midrule
\multicolumn{2}{l}{\textit{Network and API Configuration}} \\
Internet Connection & Gigabit fiber (1000 Mbps down, 500 Mbps up) \\
API Region & \texttt{us-east-1} (all providers configured for US East) \\
DNS Resolver & Cloudflare 1.1.1.1 (primary), Google 8.8.8.8 (fallback) \\
Firewall & Standard Windows Firewall (no egress restrictions) \\
\midrule
\multicolumn{2}{l}{\textit{Execution Timing}} \\
Experiment Start (UTC) & 2025-10-15 08:34:21 (timestamp of first API call) \\
Experiment End (UTC) & 2025-10-15 11:47:53 (timestamp of final CSV flush) \\
Total Elapsed Time & 3 hours, 13 minutes, 32 seconds \\
\midrule
\multicolumn{2}{l}{\textit{Randomization and Sampling}} \\
Random Seed (Global) & 42 (used for prompt shuffling and stratified sampling) \\
NumPy Random State & \texttt{np.random.seed(42)} \\
Python \texttt{random} Seed & \texttt{random.seed(42)} \\
\bottomrule
\end{tabularx}
\end{table}

\subsubsection{Model Versions and API Endpoints}

\cref{tab:model_versions_endpoints} documents the exact model identifiers and API endpoints active during the benchmark runs. Provider APIs evolve over time; these snapshots enable future researchers to request access to legacy model versions if needed for exact replication.

\begin{table}[h!]
\centering
\caption{Model versions and API endpoint URIs used in final benchmark runs.\label{tab:model_versions_endpoints}}
\begin{tabularx}{\textwidth}{@{}l l X@{}}
\toprule
\textbf{Provider} & \textbf{Model Identifier} & \textbf{API Endpoint} \\
\midrule
OpenAI & \texttt{gpt-4o-mini-2024-07-18} & \texttt{https://api.openai.com/v1/chat/completions} \\
Anthropic & \texttt{claude-3-sonnet-20240229} & \texttt{https://api.anthropic.com/v1/messages} \\
Google & \texttt{gemini-2.5-flash-latest} & \texttt{https://generativelanguage.googleapis.com/v1beta/...} \\
Cohere & \texttt{command-r-08-2024} & \texttt{https://api.cohere.ai/v1/chat} \\
Hugging Face & \texttt{meta-llama/Llama-4-Maverick-17B-128E-Instruct} & \texttt{https://api-inference.huggingface.co/models/...} \\
\bottomrule
\end{tabularx}
\end{table}

\paragraph{Model availability notes.} At the time of publication, all models remain accessible via their respective APIs under standard commercial terms. OpenAI's \texttt{gpt-4o-mini-2024-07-18} snapshot is pinned and stable; Anthropic's Claude 3 Sonnet (\texttt{20240229}) is a fixed release. Google's \texttt{gemini-2.5-flash-latest} is a rolling identifier that resolves to the current stable version; for exact replication, use the versioned identifier \texttt{gemini-2.5-flash-20240701} if still available. Cohere's Command R 08-2024 is versioned. Llama 4 Maverick is hosted on Hugging Face Inference API with model ID stability guaranteed by the repository versioning system.

\subsubsection{Cost and Pricing Metadata}

\cref{tab:pricing_metadata} records the per-million-token pricing used for cost calculations in \cref{sec:cost_modelling}. These rates were retrieved from official provider pricing pages on October 14, 2025, one day before the final benchmark runs.

\begin{table}[h!]
\centering
\caption{Pricing rates (USD per million tokens) used for cost calculations.\label{tab:pricing_metadata}}
\begin{tabularx}{\textwidth}{@{}l r r X@{}}
\toprule
\textbf{Model} & \textbf{Input (\$/1M)} & \textbf{Output (\$/1M)} & \textbf{Source URL} \\
\midrule
GPT-4o Mini & 0.15 & 0.60 & \texttt{openai.com/pricing} \\
Claude 3 Sonnet & 3.00 & 15.00 & \texttt{anthropic.com/pricing} \\
Gemini 2.5 Flash & 0.075 & 0.30 & \texttt{ai.google.dev/pricing} \\
Command R 08-2024 & 0.50 & 1.50 & \texttt{cohere.com/pricing} \\
Llama 4 Maverick & 23.50 & 23.50 & Hugging Face Inference (converted from per-second billing; see \cref{subsec:total_cost_formula}) \\
\bottomrule
\end{tabularx}
\end{table}

\paragraph{Llama pricing conversion.} Hugging Face Inference API bills Llama 4 Maverick by inference-second rather than per-token. To enable consistent cross-model cost comparisons in \cref{sec:cost_modelling,sec:results}, we converted the per-second rate (\$0.012 per inference-second) to an equivalent per-million-token price using the observed median latency (706 ms) and token counts (180 tokens combined) from our benchmark runs. The calculation yields \$47 per million tokens combined, which apportions to \$23.50/1M tokens for input and \$23.50/1M tokens for output. This pricing reflects the effective cost for our observed usage pattern and enables apples-to-apples comparisons with token-priced models. See \cref{subsec:total_cost_formula} for the detailed conversion methodology.

\paragraph{Pricing stability and future adjustments.} Providers may revise pricing over time. For longitudinal comparisons or cost-sensitivity analyses, researchers should adjust costs using updated rate tables and document the revision date. The harness logs the active price table with each run (\texttt{run\_id} metadata in SQLite; see \cref{app:sqlite_schema}), enabling retrospective cost recalculation without rerunning API calls.

\subsubsection{Artifact Provenance}

All experimental artifacts are archived in the repository under \texttt{results/} and are also available as downloadable GitHub Actions artifacts (\cref{app:github_actions_workflow}). \cref{tab:artifact_provenance} maps artifact types to storage locations and provides checksums for integrity verification.

\begin{table}[h!]
\centering
\caption{Provenance and checksums for key experimental artifacts.\label{tab:artifact_provenance}}
\begin{tabularx}{\textwidth}{@{}l X l@{}}
\toprule
\textbf{Artifact Type} & \textbf{File Path (Relative to Repository Root)} & \textbf{SHA-256 Checksum} \\
\midrule
Raw CSV Output & \texttt{results/raw\_output/run\_20251015\_0834.csv} & \texttt{5a7c9f2e...} \\
SQLite Database & \texttt{results/benchmark.db} & \texttt{d4e8b1a3...} \\
Summary JSON & \texttt{results/summary\_20251015\_0834.json} & \texttt{1f3e7c9b...} \\
Harness Logs & \texttt{logs/harness\_20251015\_0834.log} & \texttt{8c4d2f1a...} \\
Analysis Scripts & \texttt{scripts/generate\_report.py}, \texttt{statistical\_analysis.py} & \texttt{2b9e5f7c...} \\
Configuration Files & \texttt{models.json}, \texttt{system\_prompts.json}, \texttt{test\_prompts.json} & \texttt{3d4f8a1c...} \\
\bottomrule
\end{tabularx}
\end{table}

\subsubsection{Replication Instructions}

To replicate the benchmark results independently:

\begin{enumerate}
\item \textbf{Clone the repository}:
\begin{lstlisting}[language=bash,style=mystyle]
git clone https://github.com/ksawesome/llm-harness
cd llm-harness
git checkout 7ec2e57
\end{lstlisting}

\item \textbf{Set up Python environment}:
\begin{lstlisting}[language=bash,style=mystyle]
conda create -n llm-benchmark python=3.10
conda activate llm-benchmark
pip install --upgrade pip
pip install -r requirements.txt
\end{lstlisting}

\item \textbf{Configure API credentials}: Create a \texttt{.env} file in the project root with your API keys:
\begin{lstlisting}[language=bash,style=mystyle]
OPENAI_API_KEY=sk-...
ANTHROPIC_API_KEY=sk-ant-...
GOOGLE_API_KEY=AIza...
COHERE_API_KEY=...
HUGGINGFACE_API_KEY=hf_...
\end{lstlisting}

\item \textbf{Validate configuration}:
\begin{lstlisting}[language=bash,style=mystyle]
python scripts/validate_config.py
\end{lstlisting}

\item \textbf{Run the benchmark}: Execute the harness with the same parameters used in the study:
\begin{lstlisting}[language=bash,style=mystyle]
python main.py --seed 42 --concurrency 4 --verbose
\end{lstlisting}

\item \textbf{Generate analysis reports}:
\begin{lstlisting}[language=bash,style=mystyle]
python scripts/generate_report.py --input results/benchmark.db
\end{lstlisting}

\item \textbf{Verify checksums}: Compare checksums of generated artifacts against \cref{tab:artifact_provenance} to confirm bitwise reproducibility (accounting for timestamp fields).
\end{enumerate}

\paragraph{Expected deviations.} Exact bitwise reproducibility of model responses is not guaranteed due to provider-side non-determinism (temperature sampling, load balancing, infrastructure changes). However, aggregate statistics (mean latency, cost distributions, rubric scores) should remain stable within confidence intervals reported in \cref{sec:results}. Timestamp and \texttt{run\_id} fields will differ; all other telemetry fields should match closely.
